\section{Gestió dels Interessats}

% 1. VOTANTS
\begin{table}[H]
\centering
\renewcommand{\arraystretch}{1.5}
\begin{tabularx}{\textwidth}{|l|X|}
\hline
\multicolumn{2}{|c|}{\textbf{Votants (Ciutadans del cens)}} \\ \hline
\textbf{Rol} & Usuaris finals i validadors efímers de la xarxa. \\ \hline
\textbf{Perfil professional} & Públic general, persones físiques amb dret a vot actiu. \\ \hline
\multicolumn{2}{|c|}{\textbf{Avaluació i Impacte}} \\ \hline
\textbf{Interessos} & Tenen l'interès de poder exercir el seu dret a vot de manera accessible, ràpida i sense coaccions. Exigeixen privadesa total (que ningú sàpiga què han votat) i garanties que el seu vot s'ha comptat correctament. \\ \hline
\textbf{Prioritat} & \(\bullet \bullet \bullet \bullet \bullet\) \\ \hline
\textbf{Participació} & \(\bullet \bullet \bullet \bullet \bullet\) \\ \hline
\textbf{Influència} & \(\bullet \bullet \bullet \bullet \circ\) \\ \hline
\textbf{Impacte potencial} & Desús per no tenir el coneixement tècnic per llançar la instància web al seu propi sistema, si ho perceben com a insegur o massa complex, podria suposar l'abstenció deslegitima totalment el projecte. \\ \hline
\textbf{Beneficis aportats} & \textbf{Qualitatiu:} Confiança matemàtica absoluta en el procés electoral, sense dependre de la paraula de cap polític. \\ \hline
\end{tabularx}
\end{table}

%TODO: REVISAR AIXO, la xarxa ja no cau sense ells, plus, no se pot desde el movil.

\vspace{0.5cm}

% 2. CANDIDATS
\begin{table}[H]
\centering
\renewcommand{\arraystretch}{1.5}
\begin{tabularx}{\textwidth}{|l|X|}
\hline
\multicolumn{2}{|c|}{\textbf{Candidats (Partits o Opcions electorals)}} \\ \hline
\textbf{Rol} & Opcions elegibles i proveïdors d'infraestructura (Nodes Llavor). \\ \hline
\textbf{Perfil professional} & Polítics, representants legals o entitats que es presenten a les eleccions. \\ \hline
\multicolumn{2}{|c|}{\textbf{Avaluació i Impacte}} \\ \hline
\textbf{Interessos} & Necessiten garantir que les regles d'inscripció són justes i no hi ha censura prèvia. Volen que el recompte sigui completament transparent per evitar manipulacions per part dels seus rivals polítics. \\ \hline
\textbf{Prioritat} & \(\bullet \bullet \bullet \bullet \circ\) \\ \hline
\textbf{Participació} & \(\bullet \bullet \bullet \bullet \bullet\) \\ \hline
\textbf{Influència} & \(\bullet \bullet \bullet \bullet \circ\) \\ \hline
\textbf{Impacte potencial} & Podrien negar-se a reconèixer els resultats si no entenen el funcionament del \textit{blockchain}, o podrien boicotejar la infraestructura apagant els seus nodes llavor. \\ \hline
\textbf{Beneficis aportats} & \textbf{Quantitatiu i Qualitatiu:} Estalvi econòmic massiu al no haver d'imprimir paperetes ni contractar interventors o apoderats per vigilar els col·legis electorals físics. Recompte instantani. \\ \hline
\end{tabularx}
\end{table}

\vspace{0.5cm}

% 3. ORGANITZADORS
\begin{table}[H]
\centering
\renewcommand{\arraystretch}{1.5}
\begin{tabularx}{\textwidth}{|l|X|}
\hline
\multicolumn{2}{|c|}{\textbf{Organitzadors (Comitè Electoral / Estat)}} \\ \hline
\textbf{Rol} & Promotors del sistema i gestors inicials. \\ \hline
\textbf{Perfil professional} & Administració pública, ministeris o juntes electorals. \\ \hline
\multicolumn{2}{|c|}{\textbf{Avaluació i Impacte}} \\ \hline
\textbf{Interessos} & El seu objectiu és reduir els costos operatius massius d'unes eleccions tradicionals i demostrar transparència davant la comunitat internacional i els ciutadans. \\ \hline
\textbf{Prioritat} & \(\bullet \bullet \bullet \bullet \bullet\) \\ \hline
\textbf{Participació} & \(\bullet \bullet \circ \circ \circ\) \\ \hline
\textbf{Influència} & \(\bullet \bullet \bullet \bullet \bullet\) \\ \hline
\textbf{Impacte potencial} & Són els qui tenen la capacitat de finançar i adoptar el sistema legalment. Si la tecnologia els resulta aliena, poden bloquejar-ne la implantació. \\ \hline
\textbf{Beneficis aportats} & \textbf{Quantitatiu:} Reducció de milions d'euros en logística (urnes, policies, dietes, col·legis). L'ús d'una xarxa sobirana basada en Algorand elimina també les comissions per vot (\textit{fees}), fent el sistema econòmicament viable. \\ \hline
\end{tabularx}
\end{table}

\vspace{0.5cm}

% 4. AUDITORS
\begin{table}[H]
\centering
\renewcommand{\arraystretch}{1.5}
\begin{tabularx}{\textwidth}{|l|X|}
\hline
\multicolumn{2}{|c|}{\textbf{Auditors i Observadors Independents}} \\ \hline
\textbf{Rol} & Verificadors de la immutabilitat del sistema. \\ \hline
\textbf{Perfil professional} & Periodistes, observadors internacionals, ONGs i perits informàtics. \\ \hline
\multicolumn{2}{|c|}{\textbf{Avaluació i Impacte}} \\ \hline
\textbf{Interessos} & Necessiten eines que els permetin certificar que els vots no s'han manipulat durant la nit electoral, sense haver de confiar cegament en la paraula de l'organitzador. \\ \hline
\textbf{Prioritat} & \(\bullet \bullet \bullet \circ \circ\) \\ \hline
\textbf{Participació} & \(\bullet \circ \circ \circ \circ\) \\ \hline
\textbf{Influència} & \(\bullet \bullet \bullet \bullet \circ\) \\ \hline
\textbf{Impacte potencial} & La seva opinió tècnica dona validesa i prestigi al sistema de cara a l'opinió pública. Un informe negatiu d'auditoria suposa el rebuig social del projecte. \\ \hline
\textbf{Beneficis aportats} & \textbf{Qualitatiu:} Capacitat d'auditoria autònoma i asíncrona. Gràcies a l'ancoratge de l'escrutini (\textit{state anchoring}) a Ethereum, poden verificar el resultat matemàticament amb facilitat. \\ \hline
\end{tabularx}
\end{table}

\vspace{0.5cm}

% 5. ENTITATS DEMOCRÀTIQUES
\begin{table}[H]
\centering
\renewcommand{\arraystretch}{1.5}
\begin{tabularx}{\textwidth}{|l|X|}
\hline
\multicolumn{2}{|c|}{\textbf{Entitats Democràtiques i DAOs (Adoptants)}} \\ \hline
\textbf{Rol} & Reutilitzadors de la tecnologia (usuaris potencials B2B). \\ \hline
\textbf{Perfil professional} & Associacions, sindicats, universitats o Organitzacions Autònomes Descentralitzades (DAOs). \\ \hline
\multicolumn{2}{|c|}{\textbf{Avaluació i Impacte}} \\ \hline
\textbf{Interessos} & Cerquen solucions de votació segures per als seus processos interns (assemblees, eleccions sindicals, vots de socis) que no requereixin desenvolupar \textit{software} des de zero. \\ \hline
\textbf{Prioritat} & \(\bullet \bullet \circ \circ \circ\) \\ \hline
\textbf{Participació} & \(\bullet \circ \circ \circ \circ\) \\ \hline
\textbf{Influència} & \(\bullet \bullet \circ \circ \circ\) \\ \hline
\textbf{Impacte potencial} & Representen el mercat d'expansió del projecte. La seva adopció demostra que el sistema escala més enllà de l'ús governamental. \\ \hline
\textbf{Beneficis aportats} & \textbf{Quantitatiu:} Disposen d'un sistema modular de codi obert llest per ser desplegat per als seus processos interns a cost zero de desenvolupament. \\ \hline
\end{tabularx}
\end{table}

\vspace{0.5cm}

% 6. COMUNITAT DE RECERCA
\begin{table}[H]
\centering
\renewcommand{\arraystretch}{1.5}
\begin{tabularx}{\textwidth}{|l|X|}
\hline
\multicolumn{2}{|c|}{\textbf{Comunitat de Recerca i Acadèmica}} \\ \hline
\textbf{Rol} & Avaluadors tècnics i continuadors del coneixement. \\ \hline
\textbf{Perfil professional} & Professors d'universitat, investigadors en criptografia i enginyers de \textit{software}. \\ \hline
\multicolumn{2}{|c|}{\textbf{Avaluació i Impacte}} \\ \hline
\textbf{Interessos} & Interès en la innovació arquitectònica del sistema, especialment en com es resolen els problemes d'escalabilitat, la integració de la tecnologia VRF i el disseny dels "nodes efímers". \\ \hline
\textbf{Prioritat} & \(\bullet \bullet \bullet \circ \circ\) \\ \hline
\textbf{Participació} & \(\bullet \circ \circ \circ \circ\) \\ \hline
\textbf{Influència} & \(\bullet \bullet \bullet \circ \circ\) \\ \hline
\textbf{Impacte potencial} & Són els qui aproven i donen la qualificació al projecte. A llarg termini, poden obrir noves branques d'investigació a partir d'aquest prototip. \\ \hline
\textbf{Beneficis aportats} & \textbf{Qualitatiu:} Aportació d'una arquitectura sobirana ben documentada a la literatura acadèmica, facilitant una base de codi net perquè futurs estudiants puguin implementar capes de complexitat superior. \\ \hline
\end{tabularx}
\end{table}